%%%%%%%%%%%%%%%%%%%%%%%%%%%%%%%%%%%%%%%%%
% Journal Article
% LaTeX Template
% Version 2.0 (February 7, 2023)
%
% This template originates from:
% https://www.LaTeXTemplates.com
%
% Author:
% Vel (vel@latextemplates.com)
%
% License:
% CC BY-NC-SA 4.0 (https://creativecommons.org/licenses/by-nc-sa/4.0/)
%
% NOTE: The bibliography needs to be compiled using the biber engine.
%
%%%%%%%%%%%%%%%%%%%%%%%%%%%%%%%%%%%%%%%%%

%----------------------------------------------------------------------------------------
%	PACKAGES AND OTHER DOCUMENT CONFIGURATIONS
%----------------------------------------------------------------------------------------

\documentclass[
	letterpaper, % Paper size, use either a4paper or letterpaper
	10pt, % Default font size, can also use 11pt or 12pt, although this is not recommended
	unnumberedsections, % Comment to enable section numbering
	twoside, % Two side traditional mode where headers and footers change between odd and even pages, comment this option to make them fixed
]{LTJournalArticle}  

\addbibresource{patchAliasing.bib} % BibLaTeX bibliography file

\usepackage{float} % enables the [H] float specifier to pin figures/tables in place

\runninghead{Changelog} % A shortened article title to appear in the running head, leave this command empty for no running head

% \footertext{\textit{Journal of ...}} % Text to appear in the footer, leave this command empty for no footer text

\setcounter{page}{1} % The page number of the first page, set this to a higher number if the article is to be part of an issue or larger work

%----------------------------------------------------------------------------------------
%	TITLE SECTION
%----------------------------------------------------------------------------------------

\title{Changelog} % Article title, use manual lines breaks (\\) to beautify the layout

% Authors are listed in a comma-separated list with superscript numbers indicating affiliations
% \thanks{} is used for any text that should be placed in a footnote on the first page, such as the corresponding author's email, journal acceptance dates, a copyright/license notice, keywords, etc
\author{
	Matteo Boniotti, Gianluca Brignoli and Federico Sabbadini\thanks{This work is licensed under the Creative Commons Attribution 4.0 International License. Concurrently, all associated source code, probing scripts, and software artifacts generated throughout the project lifecycle are released under the open-source MIT license. \newline Copyright for any third-party material included in this work remains with its respective owners and must be respected accordingly. \newline \newline During the preparation of this work, the authors used AI-based tools, including GPT-5.5, Gemini 3.6 Flash, Sonnet 5, and Opus 4.6, to support language refinement and improve clarity. All generated or suggested content was carefully reviewed, revised where necessary, and validated by the authors, who assume full responsibility for the final content of this work.}
}

% Affiliations are output in the \date{} command
% \date{\footnotesize\textsuperscript{\textbf{1}}School of Chemistry, The University of Michigan\\ \textsuperscript{\textbf{2}}Physics Department, The University of Wisconsin\\ \textsuperscript{\textbf{3}}Biological Sciences Department, The University of Minnesota}

% Full-width abstract
%\renewcommand{\maketitlehookd}{%
	%\begin{abstract}
		%\noindent Lorem ipsum dolor sit amet, consectetur adipiscing elit. Praesent porttitor arcu luctus, imperdiet urna iaculis, mattis eros. Pellentesque iaculis odio vel nisl ullamcorper, nec faucibus ipsum molestie. Sed dictum nisl non aliquet porttitor. Etiam vulputate arcu dignissim, finibus sem et, viverra nisl. Aenean luctus congue massa, ut laoreet metus ornare in. Nunc fermentum nisi imperdiet lectus tincidunt vestibulum at ac elit. Nulla mattis nisl eu malesuada suscipit. Aliquam arcu turpis, ultrices sed luctus ac, vehicula id metus. Morbi eu feugiat velit, et tempus augue. Proin ac mattis tortor. Donec tincidunt, ante rhoncus luctus semper, arcu lorem lobortis justo, nec convallis ante quam quis lectus. Aenean tincidunt sodales massa, et hendrerit tellus mattis ac. Sed non pretium nibh. Donec cursus maximus luctus. Vivamus lobortis eros et massa porta porttitor.
	%\end{abstract}
%}

%----------------------------------------------------------------------------------------
\renewcommand{\arraystretch}{1.75}

\begin{document}
 
\maketitle % Output the title section
 
\begin{table}[!b]
\centering
\small
\renewcommand{\arraystretch}{1.4}
\begin{tabular}{@{} m{4.2cm} m{7.6cm} m{2.8cm} @{}}
\toprule
\textbf{Change} & \textbf{What was done} & \textbf{Where} \\
\midrule
Operational definition of structural aliasing
  & The representation-distance criterion $\lVert R_l(x)-R_l(y)\rVert_2\le\varepsilon$ was
    withdrawn and replaced by a readability criterion, since the probing evidence contradicts
    the near-collision reading.
  & Mathematical Formulation \\
\hline
Bayesian analysis extended to H2 and H3
  & One model per hypothesis. The H1 framework is unchanged; a phase model and two
    site-geometry models were added, with priors and decision rules.
  & Bayesian Analysis, Eqs.~(11)--(13) \\
\hline
Local-contrast hierarchy stated
  & The grouping implied by $\mu_i$ is written out: configuration, harmonic and background
    levels, plus an overlap slope $\delta_O$.
  & Bayesian Analysis, Eq.~(9) \\
\hline
Codelength response defined
  & The Gamma response is specified as a frequency-local codelength, so that the
    $\text{IsLocked}$ indicator has something to attach to.
  & Bayesian Analysis, Eq.~(10) \\
\hline
Inference protocol stated
  & Sampler settings, acceptance thresholds and the validation sequence replace the former
    one-line statement.
  & Bayesian Analysis \\
\bottomrule
\end{tabular}
\caption{Changes from the submitted Deliverable~1 to this revision.}
\label{tab:changelog}
\end{table}

This document records the changes made to Deliverable~1 \emph{after} its submitted version. The
revisions to that submission, made in response to the feedback of 4~June~2026, are recorded in the
changelog that accompanied it and are not repeated here. Only one file changed: the body section.
The equation count rises from ten to thirteen, so the contrast is now Eq.~(8) rather than Eq.~(9)
while the Gamma regression remains Eq.~(10).

\subsection{Operational definition of structural aliasing}
The submitted version defined structural aliasing by a representation-distance threshold: two
frequency-matched signals were said to alias at layer $l$ when
$\lVert R_l(x)-R_l(y)\rVert_2\le\varepsilon$ for a pre-registered $\varepsilon$. That equation has
been removed. The project's own probing does not support it, at a lock the token is a geometric
\emph{outlier} rather than a near-collision, so the criterion would have committed the report to
a claim the measurements contradict. Structural aliasing is now defined by readability: the
phenomenon is present when the signal's frequency can no longer be recovered from the model's
internal state at a locked frequency as well as at neighbouring, non-locked ones. This is the
criterion the experiments actually test, and it is what the $H1$ metrics already measure.

\subsection{Bayesian analysis: extension to H2 and H3}
The submitted version applied its two-part framework to $H1$ alone; the lecturers asked for the
treatment to be extended. Each hypothesis now has its own estimand and model, and the $H1$ models
are unchanged in substance.

Two additions are new. \emph{Phase-Invariance Analysis} ($H2$) retains the phase of each
measurement, which the contrast of Eq.~(8) averages away, cuts the phase circle into eight bins and
gives each bin its own offset (Eq.~11); the spread $\sigma_\phi$ of those offsets is tested against
a preregistered region of practical equivalence and against the phase-free model by leave-one-out.
\emph{Site-Geometry Analysis} ($H3$) targets the location of the effect rather than its size: the
token-collapse profile is regressed on indicators of the two predicted grids (Eq.~12), swept on a
grid that unions every configuration's predicted sites so no geometry is scored on a grid of its
own, with the stride-only, patch-only and null variants compared by leave-one-out; a second model
regresses the measured comb spacing on the predicted spacings $f_s/S$ and $f_s/P$ (Eq.~13), testing
the slopes $\kappa_S=1$ and $\kappa_P=0$. Because the design varies $P$ only along $P=S$,
$\kappa_P$ is declared partially confounded and the fixed-$P{=}16$ series carries the
identification.

Both models were chosen for defensibility as much as for fit: a variance component and a linear
regression on predicted values assume less, and are easier to state, than the first-harmonic and
log-log formulations they replaced.

\subsection{Local-contrast hierarchy and codelength response}
Two specifications that the submitted version left implicit are now written out. The contrast
likelihood is given in hierarchical form (Eq.~9), with configuration, harmonic and background
levels and an overlap slope $\delta_O$, and the inference is declared conditional on the single
available training seed.

The Gamma regression keeps its form and priors, but its response is now defined. The seven band
tasks span the whole operational band, so a single codelength per probe point summarises hundreds
of frequencies and leaves the $\text{IsLocked}$ indicator of Eq.~(10) nothing to attach to. The
response is therefore a \emph{frequency-local} codelength: at each candidate frequency a balanced
probe separates tones at $f_c\pm1\,\mathrm{Hz}$ from the same $[\mathrm{REG}]$ representation under
an otherwise unchanged prequential protocol. The band-task values are retained as a descriptive
cross-check.

\subsection{Inference}
The former single sentence is replaced by the full protocol: the quantities reported for each
model, the sampler settings and the acceptance thresholds ($\hat R<1.01$, effective sample sizes
above $1000$, no divergent transitions), and the validation sequence, prior predictive checks
before data collection, simulation-based parameter recovery, stratified posterior predictive
checks, and sensitivity over the prior-scale ladder.

\end{document}
